\chapter{Calcul intégral, encadrements et suites d’intégrales}
\section*{Exercice 20.1 -- Intégration par parties}
Avec $u=x$, $v'=e^x$, $u'=1$, $v=e^x$, l'intégration par parties donne $[xe^x]_0^1-\int_0^1e^x\,dx=e-(e-1)=1$. Les fonctions choisies sont continûment dérivables sur le segment.

\section*{Exercice 20.2 -- Aire entre deux courbes}
Sur $[0,1]$, $x-x^2=x(1-x)\ge0$. L'aire vaut donc $\int_0^1(x-x^2)\,dx=1/2-1/3=1/6$ unité d'aire. Inverser les fonctions aurait donné l'opposé de l'aire.

\section*{Exercice 20.3 -- Suite d’intégrales}
$0\le I_n\le\int_0^1x^n\,dx=1/(n+1)\to0$. Une intégration par parties avec $u=x^n$, $v=-e^{-x}$ donne $I_n=[-x^ne^{-x}]_0^1+n\int_0^1x^{n-1}e^{-x}\,dx$, d'où la récurrence. La borne prouve la limite sans itérer une formule numériquement instable.

\section*{Exercice 20.4 -- Somme de Riemann}
C'est une somme de rectangles à droite pour $f(x)=1/(1+x)$ sur $[0,1]$. La continuité garantit la convergence vers $\int_0^1f(x)\,dx=[\ln(1+x)]_0^1=\ln2$. Les rectangles à gauche donnent un encadrement dont l'écart vaut $1/(2n)$.

\section*{Exercice 20.5 -- Trapèzes et convexité}
La courbe est sous la corde : $f(x)\le f(a)+(x-a)(f(b)-f(a))/(b-a)$. Intégrer cette fonction affine donne exactement l'aire algébrique du trapèze, $(b-a)(f(a)+f(b))/2$. L'argument reste valable si certaines valeurs sont négatives.

\section*{Exercice 20.6 -- Série harmonique}
Sur $[k,k+1]$, $1/x\le1/k$. Donc $\ln(n+1)=\int_1^{n+1}dx/x\le H_n$. Le membre de gauche tend vers $+\infty$, donc $H_n$ aussi. Pourtant le terme $1/n$ tend vers zéro : cette condition ne suffit pas à la convergence d'une série.
